\section{Polynomial tests for bounded holomorphic functions}
\label{sec:analytic}

We will force the $w$-derivative of every bounded holomorphic function
$f$ on the tube to vanish. The elementary estimate comes from a holomorphic
disk through a point $(s_0,\zeta,w_0)$:
\[
 t\longmapsto(s_0,\zeta+t,w_0+i\alpha b t),\qquad |t|<L,
 \qquad \alpha,L>0,\quad |b|=1.
\]
If this disk lies in the tube and $|f|\le1$, Cauchy's inequality gives
\[
 |f_y(s_0,\zeta,w_0)+i\alpha b f_w(s_0,\zeta,w_0)|\le L^{-1}.
\]
We must prevent the two derivatives from canceling at many spatial
points. We arrange those points in root-of-unity cycles and hold $b$
constant on each cycle. A finite Fourier average then keeps only Taylor
degrees in one residue class. The phases will separate the prescribed
lowest-degree part of $f_w$ from the first few surviving terms of $f_y$,
with one phase choice valid for every polynomial of that degree.

The grids constructed below provide equidistributed projective directions
for polynomial sampling and cyclic affine lifts for this extraction.
Their additional separation and shell estimates will be used in
Section~\ref{sec:assembly} to localize the shear coefficients in disjoint
spatial neighborhoods when realizing the disks. After constructing the grids
and choosing the phases, we formulate the precise disk conditions and
prove that repeated tests force the holomorphic Liouville property.

Throughout this section, $\sigma$ denotes the unitary-invariant
probability measure on $\mathbb P^{n-1}$.  All norms on vectors and
derivatives are Euclidean unless indicated otherwise.  Constants may
depend on the fixed dimension $n$; the exponents in the separation
estimate below do not.

\subsection{Polynomial grids and their affine lifts}

\begin{lemma}[Homogeneous polynomial grids]
\label{grid:polynomials}
Fix $n\ge2$.  For every sufficiently large integer $N$ one can choose a
system of homogeneous polynomials of degree $N$,
\[
 Q_N=(H_N,p_N):\mathbb C^n\longrightarrow\mathbb C^{n-1}\times\mathbb C,
\]
with the following properties.  These choices can be made along all
integers $N\to\infty$, or along any prescribed sequence of degrees
tending to infinity.
\begin{enumerate}
\item There is a constant $C_n$ such that
\begin{equation}
 \label{grid:sphere-bounds}
 N^{-C_n}\le |Q_N(u)|\le N\qquad (|u|=1).
\end{equation}
\item The projective zero set of $H_N$ consists of
$M_N=N^{n-1}$ transverse points.  A subset $\mathcal G_N$ of these
points satisfies
\[
 |\mathcal G_N|=(1-o(1))M_N,
 \qquad
 \frac1{M_N}\sum_{\xi\in\mathcal G_N}\delta_\xi
 \ \rightharpoonup\ \sigma.
\]
For every $\xi\in\mathcal G_N$ and every unit representative $u$ of
$\xi$,
\begin{equation}
 \label{grid:goodness}
 N^{-1}\le |p_N(u)|\le N,
 \qquad
 \operatorname{sing}_{\min}
       \bigl(DH_N(u)|_{u^\perp}\bigr)\ge N^{-1}.
\end{equation}
\item Above each $\xi\in\mathcal G_N$, choose any lift $v_\xi$ such
that $p_N(v_\xi)=1$.  All $N$ lifts
\[
 e^{2\pi i q/N}v_\xi,\qquad 0\le q<N,
\]
have radii in $[N^{-1/N},N^{1/N}]$.  With the absolute choice
$k_*=30$, the Euclidean balls of radius $2N^{-k_*}$ about all these
lifts are pairwise disjoint.  If $v$ is one of the lifts, then
\begin{equation}
 \label{grid:shell-bound}
 |(H_N(z),p_N(z)-1)|\ge N^{-k_*-10}
 \quad\hbox{when}\quad
 N^{-k_*}\le |z-v|\le2N^{-k_*}.
\end{equation}
In a neighborhood of radius $4N^{-k_*}$ of every such lift, the first
and second derivatives of $Q_N$ have norm $O_n(N^5)$.
\end{enumerate}
In particular, weak approximation to any prescribed finite collection
of continuous test functions, with any fixed positive accuracy, can be
required by taking $N$ sufficiently large.
\end{lemma}

\begin{proof}
We use the standard unitary-invariant Gaussian ensemble of homogeneous
polynomials; its covariance and conditional zero-density framework appear
in \citet[Sections~2.2 and~2.4]{BleherShiffmanZelditch2000}. The proof
below supplies the needed formulas and all additional quantitative grid
properties. Randomness is used only to establish existence. A
standard circular complex Gaussian has density
$\pi^{-1}e^{-|z|^2}$ on $\mathbb C$.  Choose the components of $Q_N$
independently, with
\[
 Q_N^{(a)}(z)=
 \sum_{|\beta|=N}\left(\frac{N!}{\beta!}\right)^{1/2}
       g_{a,\beta}z^\beta,
 \qquad
 \mathbb E\bigl[Q_N^{(a)}(z)
       \overline{Q_N^{(a)}(u)}\bigr]=(z\cdot\bar u)^N.
\]
Here all $g_{a,\beta}$ are independent standard circular complex
Gaussians.  The covariance identity is the multinomial formula.  In
particular, at a unit vector the components of $Q_N$ are independent
standard complex Gaussians, and the law is invariant under unitary
changes of coordinates.

We first prove \eqref{grid:sphere-bounds} with probability tending to
one.  If $A=\sup_{|z|=1}|Q_N(z)|$, homogeneity gives
$\sup_{|z|\le1+1/N}|Q_N(z)|\le eA$.  Cauchy's inequality on complex
lines therefore bounds the first derivative on the unit sphere by
$CNA$, with an absolute constant $C$.  A sphere net of mesh $c/N$,
with a sufficiently small absolute $c$, consequently contains a point
where the norm is at least $A/2$.  Such a net has
$O_n(N^{2n-1})$ points.  Gaussian tails show
\[
 \mathbb P(A>N)\le C_n' N^{2n-1}e^{-c_n' N^2}\longrightarrow0.
\]
On the event $A\le N$, the Lipschitz constant on the sphere is at
most $CN^2$.  Use a second net of mesh $cN^{-C_n-2}$.  If the norm
somewhere on the sphere is less than $N^{-C_n}$, it is less than
$2N^{-C_n}$ at a point of this net.  A standard complex Gaussian
vector in $\mathbb C^n$ has probability at most $C_n' r^{2n}$ of
lying in a ball of radius $r$.  The union bound is therefore at most
\[
 C_n' N^{(C_n+2)(2n-1)-2nC_n}
   =C_n' N^{-C_n+4n-2}.
\]
Choose, for example, $C_n=8n$.  This proves both sphere bounds with
probability tending to one.

Put $d=n-1$.  Almost surely the projective zeros of the $d$-component
system $H_N$ are transverse.  To see this directly, let $\mathcal A$
be its finite-dimensional complex coefficient space.  Evaluation at
any point of $\mathbb P^d$ maps $\mathcal A$ onto the fiber of
$\mathcal O(N)^{\oplus d}$.  Hence the universal zero incidence
manifold is smooth, and a coefficient vector is a regular value of its
projection to $\mathcal A$ precisely when the corresponding section
has only transverse zeros.  Sard's theorem gives a Lebesgue-null
exceptional set; the Gaussian law is absolutely continuous.
Transversality and compactness imply that the zero set is finite.
The projective B\'ezout theorem for a zero-dimensional intersection of
$d$ hypersurfaces of degrees $N,\ldots,N$ gives total intersection
number $N^d$; every transverse intersection has multiplicity one
\citep[Theorem~III-71]{EisenbudHarris2000GeometrySchemes}.  Thus the number of zeros is the deterministic
quantity $M_N=N^d$ almost surely.

We next justify the counting formulas needed for equidistribution.
They are the one- and two-point, zero-dimensional cases of the
correlation formula in \citet[Theorem~2.1]{BleherShiffmanZelditch2000}.
Our subsequent off-diagonal argument proves the particular convergence
needed here; the more general simultaneous-zero convergence theorem of
\citet[Corollary~1.3]{Shiffman2008Convergence} supplies a broader context.
In a holomorphic coordinate chart of $\mathbb P^d$, regard the system
as a holomorphic map $H_a(z)$ depending linearly on its coefficients.
Locally split those coefficients as $(a,b)$ so that
\[
 H_{a,b}(z)=E(z)a+R(z,b),\qquad \det E(z)\ne0.
\]
This is possible because evaluation is surjective.  The equation
$H_{a,b}(z)=0$ is equivalent to
$a=A(z,b):=-E(z)^{-1}R(z,b)$.  At a zero, with
$J=D_zH_{a,b}(z)$,
\[
 D_zA=-E(z)^{-1}J.
\]
For fixed $b$, the change-of-variables formula with multiplicity for
$z\mapsto A(z,b)$ thus introduces the real Jacobian
$|\det E|^{-2}|\det J|^2$.  Integrating in $b$ against the smooth
Gaussian coefficient density yields the one-point zero density
\begin{equation}
 \label{grid:one-point-density}
 \rho_{1,N}(z)=
 p_{H_N(z)}(0)\,
 \mathbb E\bigl[|\det DH_N(z)|^2\mid H_N(z)=0\bigr].
\end{equation}
The same argument works with a bounded measurable mark multiplying
the zero count.  It first applies to nonnegative marks, and the
general bounded case follows by decomposition.

For two distinct projective points, joint evaluation onto the two
fibers is surjective.  Indeed, powers of linear forms that vanish at
one point and not the other give independent prescribed values for
each component.  Solve now for $2d$ coefficient variables and apply
the same change of variables in the two point coordinates.  The
point Jacobian is block diagonal, giving
\begin{equation}
 \label{grid:two-point-density}
 \begin{split}
 \rho_{2,N}(z,z')={}&p_{(H_N(z),H_N(z'))}(0,0)\\
 &\quad\cdot
 \mathbb E\bigl[
   |\det DH_N(z)|^2|\det DH_N(z')|^2
       \mid H_N(z)=H_N(z')=0\bigr].
 \end{split}
\end{equation}
Partitions of unity give these formulas on projective space and on
the complement of its diagonal, respectively.  This derivation uses
only the coefficient density and ordinary change of variables; no
independence of zeros has been assumed.

At a point represented by a unit vector $u$, take affine tangent
coordinates $u+\sum_{j=1}^d z_j e_j$ with $(e_j)$ orthonormal in
$u^\perp$.  Differentiating the covariance shows that $H_N(u)$ is
independent of $N^{-1/2}DH_N(u)|_{u^\perp}$, and that the latter is
a $d\times d$ matrix of independent standard complex Gaussians.
At two distinct projective points, the cross-covariances of these
normalized value and first-derivative data tend to zero locally
uniformly off the diagonal.  More explicitly, on a compact set of
pairs with $|u\cdot\bar v|\le\rho<1$, they are bounded by a fixed
power of $N$ times $\rho^{N-2}$.  This follows by differentiating
$(u\cdot\bar v)^N$ at most once at each point.  The joint Gaussian
covariance consequently converges uniformly there to the direct sum
of its one-point covariances.  Conditional Gaussian covariances given
the values zero converge as well.  Gaussian moments of the fixed
polynomials $|\det J|^2$ and their products therefore converge.
Equations \eqref{grid:one-point-density} and
\eqref{grid:two-point-density} imply that the two-point density,
divided by $M_N^2$, converges locally uniformly off the diagonal to
the product of the normalized one-point densities.

The exact total number of zeros lets us pass from off-diagonal
decorrelation to weak convergence without estimating the two-point
density near the diagonal. Let
\[
 \mu_N=\frac1{M_N}\sum_{H_N(\xi)=0}\delta_\xi.
\]
Unitary invariance and the deterministic count imply
$\mathbb E\mu_N=\sigma$.  The positive measure
$\mathbb E(\mu_N\otimes\mu_N)$ has total mass one.  Off the
diagonal, its density converges to that of $\sigma\otimes\sigma$,
by the preceding calculation.  This implies convergence on the
whole product: given $\varepsilon>0$, choose a compact subset of the
off-diagonal region having $(\sigma\otimes\sigma)$-mass greater
than $1-\varepsilon$.  Its mass under
$\mathbb E(\mu_N\otimes\mu_N)$ is greater than $1-2\varepsilon$
for all sufficiently large $N$.  At most $2\varepsilon$ of the
total mass can remain elsewhere.  Since the diagonal has zero
$(\sigma\otimes\sigma)$-mass, there is no additional limiting
mass there.  Thus
\[
 \mathbb E(\mu_N\otimes\mu_N)
       \rightharpoonup\sigma\otimes\sigma.
\]
For every continuous test function, its integral against $\mu_N$
therefore has variance tending to zero.  This proves weak
convergence to $\sigma$ in probability.

It remains to discard the bad zeros.  Use
\eqref{grid:one-point-density} with marks.  If $G$ is a $d\times d$
standard complex Gaussian matrix, then at a zero the derivative
condition in \eqref{grid:goodness} fails only when
$\operatorname{sing}_{\min}(G)<N^{-3/2}$.  Its expected discarded
fraction is
\[
 \frac{\mathbb E\bigl[
       |\det G|^2
       \boldsymbol{1}_{\{\operatorname{sing}_{\min}(G)<N^{-3/2}\}}
       \bigr]}{\mathbb E|\det G|^2}
       \longrightarrow0.
\]
Indeed a Gaussian square matrix is nonsingular almost surely,
$|\det G|^2$ is integrable, and dominated convergence applies.
The independent polynomial $p_N$ has a standard complex Gaussian
value at every unit representative.  Its two forbidden tails have
probability
\[
 \mathbb P(|p_N|<N^{-1})+\mathbb P(|p_N|>N)
   =1-e^{-N^{-2}}+e^{-N^2}\longrightarrow0.
\]
Markov's inequality shows that the total discarded fraction tends
to zero in probability.  Removing that fraction does not alter weak
convergence.  The modulus and singular-value conditions are
unchanged on multiplying a unit representative by a phase, so they
hold for every such representative.

We may now select deterministic systems.  Take a countable uniformly
dense family of continuous functions on projective space, and impose
successively accurate conditions on successively longer finite
subfamilies.  Convergence in probability, the vanishing bad
fraction, and the probability-one transversality statement show that
these conditions and \eqref{grid:sphere-bounds} have a nonempty
intersection for every sufficiently large degree.  Choose one system
from that intersection.  A diagonal choice gives all the asserted
limits and works equally well along any prescribed sequence of
degrees.

We have obtained the projective sampling properties. It remains to
provide disjoint neighborhoods in which Section~\ref{sec:assembly} can
localize its shear coefficients. Fix a good direction and a lift
$v$ with $p_N(v)=1$.
Homogeneity and \eqref{grid:goodness} give
$N^{-1/N}\le |v|\le N^{1/N}$.  In the splitting of the domain into
the radial complex line and its Hermitian orthogonal complement,
$DH_N(v)$ vanishes on the radial line, while
$Dp_N(v)$ has radial derivative $N/|v|$.  The transverse derivative
of $H_N$ has least singular value at least $N^{-4}$ for large $N$:
the scaling from a unit representative costs at most one additional
factor $N$, so this is a deliberately weaker bound than necessary.
Homogeneity, the sphere upper bound, and Cauchy's inequality at
scale $c/N$ give first- and second-derivative bounds $O_n(N^5)$ in
the stated neighborhood.  At the lift itself, the transverse first
derivatives are $O_n(N^4)$.  Block inversion therefore gives
\[
 \bigl\|D(H_N,p_N-1)(v)^{-1}\bigr\|\le N^9
\]
for all sufficiently large $N$.  Taylor's formula, uniformly for
$|z-v|\le4N^{-30}$, now yields
\[
 |(H_N(z),p_N(z)-1)|
 \ge N^{-9}|z-v|-C_n' N^5|z-v|^2
 \ge\tfrac12 N^{-9}|z-v|.
\]
This proves \eqref{grid:shell-bound}.  It also excludes every other
zero of $(H_N,p_N-1)$ within distance $4N^{-30}$, proving
disjointness of the balls about all good lifts.
\end{proof}

\subsection{Choosing phases uniformly for a fixed Taylor order}

The grids now supply many separated evaluation points. Their next role
is linear algebraic: we choose a phase at each projective point so that
no fixed-order polynomial jet can be hidden among the first few Taylor
degrees surviving a root-of-unity average. The independent-phase moment
calculation is the same quadratic-form identity used for random-phase
trace estimates in \citet{IitakaEbisuzaki2004}.
The additional rank and polynomial superlevel estimates give one phase
choice uniformly for an entire coefficient space.

\begin{lemma}[Phases separate low-degree jets]
\label{grid:phases}
Fix positive integers $k,n$, with $n\ge2$, such that
\begin{equation}
 \label{grid:dimension-choice}
 \sum_{j=1}^{k-1}\frac{j^{n-1}}{(n-1)!}<\frac1{100}.
\end{equation}
Fix a nonnegative integer $h$.  For sufficiently large polynomial
grids from Lemma~\ref{grid:polynomials}, choose one lift $v_\xi$
above each good direction and let $V_N\subset\mathbb C^{\mathcal G_N}$
be the space of evaluation vectors of sums of homogeneous
polynomials of degrees
\[
 h,\ h+N,\ \ldots,\ h+(k-1)N.
\]
Let $\Pi_N$ be orthogonal projection onto $V_N$ for the unweighted
Euclidean point norm.  There are phases $b_\xi$, $|b_\xi|=1$, and a
constant $\delta_h>0$, independent of $N$ and of the chosen lift
phases, such that
\begin{equation}
 \label{grid:phase-lower-bound}
 M_N^{-1/2}
 \left\|(1-\Pi_N)
       \bigl(b_\xi P(v_\xi)\bigr)_{\xi\in\mathcal G_N}\right\|
 \ge\delta_h\|P\|
\end{equation}
for every homogeneous polynomial $P$ of degree $h$.  Here $\|P\|$
is any fixed Euclidean norm on its coefficient space, chosen
independently of the grid.  All required grid accuracies and degree
thresholds depend only on the fixed $n,k,h$ and this norm.
\end{lemma}

\begin{proof}
Dimension gives
\[
 \frac{\operatorname{rank}\Pi_N}{M_N}
 \le N^{1-n}\sum_{j=0}^{k-1}
       \binom{h+jN+n-1}{n-1}<\frac2{100}
\]
for all sufficiently large $N$.  Indeed its limit is the sum in
\eqref{grid:dimension-choice}.

There is a number $t_h>0$ such that, on every sufficiently accurate
grid, every polynomial with $\|P\|=1$ satisfies
\begin{equation}
 \label{grid:uniform-superlevel}
 |P(v_\xi)|>t_h
 \quad\hbox{at at least }0.8M_N\hbox{ good directions}.
\end{equation}
For completeness, a nonzero homogeneous polynomial has a
projective zero set of $\sigma$-measure zero.  In affine charts this
is the elementary fact that the zero set of a nonzero complex
polynomial has Lebesgue measure zero, obtained by induction and
Fubini's theorem.  Thus, for each norm-one polynomial $P$, some
positive threshold has a superlevel set of measure greater than
$0.9$.  Uniform continuity in the coefficients makes the same
superlevel estimate, with a smaller threshold, valid in a
neighborhood of $P$.  The coefficient unit sphere is compact, so
finitely many such neighborhoods suffice and their thresholds have
a positive common lower bound.  Weak convergence of the grid
measures applied to these finitely many open superlevel sets gives
the estimate for unit representatives, with a fixed margin in both
threshold and counting proportion.  Since
$|v_\xi|^h\to1$ uniformly, it gives
\eqref{grid:uniform-superlevel}.  This proof requires only finitely
many fixed approximation conditions for the given $h$.

The projection has $0\le(\Pi_N)_{\xi\xi}\le1$ and diagonal sum
less than $0.02M_N$.  Consequently at most $0.04M_N$ diagonal
entries exceed $1/2$.  Choose the $b_\xi$ temporarily as independent
uniform phases.  Writing $Q_N'=1-\Pi_N$, independence gives
\[
 \begin{split}
 &\mathbb E\left[
 M_N^{-1}\left\|Q_N'
       \bigl(b_\xi P(v_\xi)\bigr)_\xi\right\|^2\right]\\
 &\hspace{25mm}=
 M_N^{-1}\sum_\xi
       \bigl(1-(\Pi_N)_{\xi\xi}\bigr)|P(v_\xi)|^2
 \ge0.38t_h^2
 \end{split}
\]
when $\|P\|=1$.  The last inequality uses the intersection of the
$0.8M_N$ indices in \eqref{grid:uniform-superlevel} with the
complement of the at most $0.04M_N$ high-leverage indices.

We need a single phase choice for all polynomials, rather than a
choice depending on $P$.  Let $(e_a)$ be a fixed orthonormal basis
of the degree-$h$ coefficient space, and put $q_a(\xi)=e_a(v_\xi)$.
These evaluations have a common bound depending only on $n,h$ and
the chosen norm.  The Gram matrix of the random linear map in
\eqref{grid:phase-lower-bound} has entries
\[
 G_{ab}=M_N^{-1}\sum_{\xi,\eta}
       \overline{q_a(\xi)}\,\overline{b_\xi}
       (Q_N')_{\xi\eta}b_\eta q_b(\eta).
\]
Its diagonal-index terms are deterministic.  For distinct indices,
the products $\overline{b_\xi}b_\eta$ are orthonormal in
$L^2$ of the phase probability space: two such ordered products
have nonzero inner product only when their ordered pairs agree.
In particular, reversal of the pair gives expectation
$\mathbb E(\overline{b_\xi}^{\,2}b_\eta^2)=0$.
It follows that
\[
 \begin{split}
 \mathbb E|G_{ab}-\mathbb EG_{ab}|^2
 &\le C_{n,h}M_N^{-2}
       \sum_{\xi\ne\eta}|(\Pi_N)_{\xi\eta}|^2\\
 &\le C_{n,h}M_N^{-2}\operatorname{rank}\Pi_N
   =O_{n,h}(M_N^{-1}).
 \end{split}
\]
The Gram matrix has fixed finite size.  Chebyshev's inequality
therefore implies that its operator-norm distance from its mean
tends to zero in probability.  Its mean is bounded below by
$0.38t_h^2$ times the identity.  A phase choice with least
eigenvalue at least $0.19t_h^2$ exists for every sufficiently large
grid.  Taking $\delta_h=\sqrt{0.19}\,t_h$ proves the assertion.
All estimates used only absolute values of lift evaluations, so
the threshold and lower bound are uniform over their arbitrary
representative phases.
\end{proof}

\begin{remark}
\label{grid:finite-dimension}
For every fixed $k$, condition \eqref{grid:dimension-choice} holds
in some sufficiently large finite dimension.  Indeed, with
$r=n-1$, its left side is bounded by
$(k-1)^{r+1}/r!$, which tends to zero as $r\to\infty$.
The argument will use one such fixed dimension, not an assertion
in every dimension.
\end{remark}

\subsection{From disk estimates to the holomorphic Liouville property}

\begin{lemma}[Root-of-unity extraction]
\label{grid:frequency}
Let $0\le h<N$, and suppose $g$ is holomorphic on a neighborhood of
the closed ball $|s|\le S$, with $|g|\le K$ there.  Write
$g=\sum_{r\ge0}g_r$ for its homogeneous Taylor expansion at zero.
If $R|v|<S$, then
\[
 \frac1N\sum_{q=0}^{N-1}e^{-2\pi i hq/N}
       g(Re^{2\pi i q/N}v)
   =\sum_{j\ge0}g_{h+jN}(Rv).
\]
If $R|v|/S\le e^{-0.99c/N}$ and $c\ge1$, truncating the sum after
$j=k-1$ has error at most $C K e^{-0.99kc}$, while retaining just
the term $g_h(Rv)$ has error at most $C K e^{-0.99c}$.  The constant
$C$ is absolute, independent of $N,h,R,S$.
\end{lemma}

\begin{proof}
Cauchy's inequality on the complex line through $v$ gives
\[
 |g_r(Rv)|\le K(R|v|/S)^r.
\]
The homogeneous expansion converges absolutely at the evaluation
points.  The finite Fourier average annihilates every degree not
congruent to $h$ modulo $N$, giving the identity.  Put
$\theta=R|v|/S$.  The first tail is bounded by
\[
 K\sum_{j=k}^{\infty}\theta^{h+jN}
   \le\frac{K e^{-0.99kc}}{1-e^{-0.99c}}.
\]
The case $k=1$ gives the other estimate.  The denominator is
bounded away from zero for $c\ge1$.
\end{proof}

We now convert these finite-dimensional estimates into the analytic
condition needed by the geometric construction. A large sheared disk
bounds a linear combination of the $y$- and $w$-derivatives. The phases
separate the degree-$h$ part of the latter from the former. Heights
$D_j$ growing more slowly than the accuracy parameters $c_j$ will let
the three-lines theorem turn smallness at those heights into vanishing.

\begin{proposition}[A testing criterion for the holomorphic Liouville property]
\label{grid:criterion}
Fix constants $p_0,p_1\ge0$ and $m_0>p_1$.  Choose an integer $k$
and a finite dimension $n\ge2$ so that
\begin{equation}
 \label{grid:exponent-choice}
 0.99k>p_0+p_1+1,
 \qquad
 \sum_{r=1}^{k-1}\frac{r^{n-1}}{(n-1)!}<0.01.
\end{equation}
Let $(\zeta_j,h_j)$ be a sequence of pairs from a countable dense
subset of an open disk in $\mathbb C$ and the nonnegative integers,
with every pair repeated infinitely often.  At stage $j$, choose
numbers $R_j\ge1$, $c_j\ge1$, an integer $N_j>h_j$, and a grid
from Lemma~\ref{grid:polynomials} sufficiently large and accurate
for Lemma~\ref{grid:phases} at order $h_j$.  Set
\[
 S_j=R_j e^{c_j/N_j},\qquad \log N_j\le c_j/100.
\]
Choose heights $D_j$ and positive numbers $\alpha_j$ such that
\[
 \alpha_j^{-1}\le e^{p_1c_j},
\]
and, along the repetitions of every fixed pair,
\begin{equation}
 \label{grid:height-rates}
 c_j\longrightarrow\infty,\qquad
 D_j\longrightarrow\infty,\qquad
 \frac{D_j}{c_j}\longrightarrow0.
\end{equation}
There are phases $b_{j,\xi}$, one per good projective direction
and constant throughout its cycle of lifts, with the following
property.

Let $\ell$ be a smooth real function on $\mathbb C^n\times\mathbb C$,
and let
\[
 \mathcal T=\{(s,y,w):\operatorname{Im}w>\ell(s,y)\}.
\]
Suppose every stage satisfies these conditions:
\begin{enumerate}
\item On $|s|\le S_j$,
\begin{equation}
 \label{grid:height-gap}
 \ell(s,\zeta_j)\le D_j-2.
\end{equation}
\item For every holomorphic $f$ on $\mathcal T$ with $|f|\le1$,
every $w$ with $\operatorname{Im}w=D_j$, and every $|s|=S_j$,
\begin{equation}
 \label{grid:sphere-derivative}
 |f_y(s,\zeta_j,w)|\le e^{p_0c_j}.
\end{equation}
\item For every good direction $\xi$, every lift
$v=e^{2\pi iq/N_j}v_\xi$, and every $w$ on that height line, the
holomorphic disk
\begin{equation}
 \label{grid:sheared-disk}
 t\longmapsto
 \bigl(R_jv,\ \zeta_j+t,\ w+i\alpha_j b_{j,\xi}t\bigr),
 \qquad |t|<e^{m_0c_j},
\end{equation}
is contained in $\mathcal T$.
\end{enumerate}
Then every bounded holomorphic function on $\mathcal T$ is constant.
The phases depend only on the grids and their assigned jet orders,
not on $\ell$, $f$, or $w$.
\end{proposition}

\begin{proof}
For each stage choose the phases supplied by
Lemma~\ref{grid:phases}.  Assign its phase to all $N_j$ lifts of
the same direction.  We prove the assertion for a holomorphic
function with $|f|\le1$, which suffices by constant rescaling.

Fix a stage and temporarily omit its subscript.  For each $w$ with
$\operatorname{Im}w=D$, put
\[
 A(s)=f_y(s,\zeta,w),\qquad I(s)=i f_w(s,\zeta,w).
\]
Both are holomorphic on a neighborhood of the closed spatial ball
$|s|\le S$, by the strict height margin and compactness.  The
maximum principle and \eqref{grid:sphere-derivative} give
$|A|\le e^{p_0c}$ on the ball.  The unit $w$-disk centered at any
$(s,\zeta,w)$ in this ball lies in $\mathcal T$, by
\eqref{grid:height-gap}; Cauchy's inequality gives $|I|\le1$.
On a sheared disk, the derivative at $t=0$ is
$A(Rv)+\alpha b_\xi I(Rv)$.  Thus
\begin{equation}
 \label{grid:node-inequality}
 |A(Rv)+\alpha b_\xi I(Rv)|\le e^{-m_0c}.
\end{equation}
Every estimate so far is uniform in the real part of $w$.

For each good direction take the Fourier average in
Lemma~\ref{grid:frequency} over its cycle.  Since $b_\xi$ is
constant on that cycle, it factors out of the averaged second
term of \eqref{grid:node-inequality}.  The lift radius bound gives
\[
 \frac{R|v_\xi|}{S}
 \le\exp\left(\frac{\log N-c}{N}\right)
 \le e^{-0.99c/N}.
\]
Let $a$ and $i'$ be the resulting vectors of averaged $A$ and $I$
values, respectively.  There is a vector $a_0\in V_N$ such that
\[
 \|a-a_0\|_\infty\le C e^{(p_0-0.99k)c}.
\]
If $P$ is the homogeneous degree-$h$ Taylor term of the function
$s\mapsto I(Rs)$, then
\[
 \bigl\|i'-(P(v_\xi))_\xi\bigr\|_\infty
       \le C e^{-0.99c}.
\]
The averaged node inequalities still have supremum norm at most
$e^{-m_0c}$.  Apply $1-\Pi_N$ and use $a_0\in V_N$.  Since the
number of good directions is at most $M_N$, conversion from the
supremum norm to $M_N^{-1/2}$ times Euclidean norm costs at most
one.  Lemma~\ref{grid:phases} consequently gives
\[
 \begin{split}
 \alpha\delta_h\|P\|
 &\le e^{-m_0c}
       +C e^{(p_0-0.99k)c}
       +C\alpha e^{-0.99c},\\
 \|P\|
 &\le C_h\left(
       e^{(p_1-m_0)c}
       +e^{(p_1+p_0-0.99k)c}
       +e^{-0.99c}\right).
 \end{split}
\]
In particular, with the fixed positive number
\[
 \varepsilon=\min\{m_0-p_1,\ 0.99k-p_0-p_1,\ 0.99\},
\]
we have
\begin{equation}
 \label{grid:jet-decay}
 \|P\|\le C_h e^{-\varepsilon c}.
\end{equation}
The constants are independent of the stage once $h$ is fixed.
This is exactly the uniformity in Lemma~\ref{grid:phases}; the
Fourier tail denominator is $1-e^{-0.99c}$, not
$1-e^{-0.99c/N}$.  Since every degree-$h$ coefficient of $I(Rs)$
is $R^h$ times the corresponding coefficient of $I(s)$ and
$R\ge1$, \eqref{grid:jet-decay} bounds each degree-$h$ coefficient
of $f_w(s,\zeta,w)$ by the same exponential rate.

We have obtained exponential decay for every coefficient of the
prescribed degree, uniformly along the entire horizontal $w$-line.
It remains to bring that estimate down from the growing testing heights
to a fixed open part of the tube.

Fix now a scheduled center $\zeta$ and a multiindex $\beta$ of
length $h$.  The coefficient
\[
 g_\beta(w)=\frac1{\beta!}
       \partial_s^\beta f_w(0,\zeta,w)
\]
is holomorphic on the whole connected half-plane
$\operatorname{Im}w>\ell(0,\zeta)$.  At every point of this
half-plane a spatial neighborhood is available, so the derivative
is well-defined and holomorphic there.  Choose once a small
closed spatial polydisk about zero and a number $t_0$ exceeding
the supremum of $\ell(s,\zeta)$ on that polydisk by more than two.
Cauchy's inequalities in $s$ and on a unit $w$-disk give
\[
 |g_\beta(w)|\le B_\beta
 \qquad (\operatorname{Im}w\ge t_0),
\]
where $B_\beta$ is independent of every testing stage and of
$\operatorname{Re}w$.  Along the stages assigned to $(\zeta,h)$,
\eqref{grid:jet-decay} gives
\[
 \sup_{\operatorname{Im}w=D_j}|g_\beta(w)|
       \le C_h e^{-\varepsilon c_j}.
\]
For any fixed $t>t_0$, apply Hadamard's three-lines theorem
\cite[Theorem~7.13]{Tropp2022} to the horizontal strip between heights
$t_0$ and $D_j$, once $D_j>t$.
Writing $\theta_j=(t-t_0)/(D_j-t_0)$, it yields
\[
 \sup_{\operatorname{Im}w=t}|g_\beta(w)|
 \le B_\beta^{\,1-\theta_j}
       (C_h e^{-\varepsilon c_j})^{\theta_j}.
\]
If $B_\beta=0$ there is nothing to prove.  Otherwise, the
logarithm of the right side is
\[
 \log B_\beta+
 \frac{t-t_0}{D_j-t_0}
    \bigl(\log(C_h/B_\beta)-\varepsilon c_j\bigr)
       \longrightarrow-\infty
\]
by \eqref{grid:height-rates}.  Hence $g_\beta$ vanishes in the
upper half-plane and, by the identity theorem, throughout
$\operatorname{Im}w>\ell(0,\zeta)$.  Repeating this argument for
each order proves that all spatial Taylor coefficients of $f_w$
at zero vanish at every scheduled center.  Dependence of the
constants and of the preliminary height $t_0$ on the order causes
no difficulty: each order has its own infinite sequence, and
analytic continuation gives the same full half-plane afterwards.

To pass from these statements to a common open set, choose a
spatial ball about zero and an open subdisk of the center disk,
both with compact closures.  Choose $L$ larger than the supremum
of $\ell$ on their product.  On the product with
$\operatorname{Im}w>L$, the function $f_w$ is jointly holomorphic.
At each scheduled center in the subdisk, all of its spatial
Taylor coefficients vanish for every such $w$.  Taylor uniqueness
and the identity theorem on the spatial ball give vanishing on
that ball.  Continuity in $y$ and density of the scheduled centers
then give vanishing on the whole product.  The tube is connected:
the coordinates
\[
 (s,y,\operatorname{Re}w,
        \operatorname{Im}w-\ell(s,y))
\]
identify it with
$\mathbb C^{n+1}\times\mathbb R\times(0,\infty)$.
The several-variable identity theorem therefore gives $f_w=0$
everywhere on $\mathcal T$.

For every fixed $(s,y)$ the $w$-fiber is a connected half-plane,
so $f$ is constant along that fiber.  It descends to a bounded
function $F$ on $\mathbb C^{n+1}$.  This descended function is
holomorphic: near any base point choose a single constant $w_0$
whose imaginary part exceeds the supremum of $\ell$ on a small
base neighborhood, and write $F(s,y)=f(s,y,w_0)$ there.  These
local holomorphic definitions agree on overlaps by fiber
constancy.  Liouville's theorem, applied on complex lines in
$\mathbb C^{n+1}$, makes $F$, and therefore $f$, constant.
\end{proof}

\begin{remark}
\label{grid:geometric-interface}
Condition \eqref{grid:sphere-derivative} follows immediately if,
at the stated sphere points and heights, the unsheared $y$-disk
of radius $e^{-p_0c_j}$ lies in the tube.  Thus the criterion asks
the geometric construction for only three quantitative facts: a
height bound on a spatial ball, sufficiently wide unsheared disks
on its boundary, and much wider disks in the prescribed sheared
directions at its polynomial nodes.  The phase choices are made
before any bounded holomorphic function is considered.
\end{remark}
